\documentclass[10pt]{article}

\usepackage{titlesec}
\usepackage{geometry}
\geometry{verbose,tmargin=.9in,bmargin=.9in,lmargin=1.0in,rmargin=1.0in}

\usepackage{amsmath,amsfonts,amssymb}
\usepackage{bbm}

\usepackage[usenames,dvipsnames,svgnames,table]{xcolor}
\usepackage[colorlinks=true, linkcolor=red, urlcolor=blue, citecolor=gray]{hyperref}
\usepackage{enumitem}

\definecolor{nyuDarkPurple}{HTML}{330662}
\definecolor{nyuOfficialPurple}{HTML}{57068c}

\newcommand{\spara}[1]{\vspace{.5em}\noindent {\large\sffamily\textcolor{nyuOfficialPurple}{#1}}}
\titleformat{\section}[hang]{\Large\sffamily\color{nyuDarkPurple}}{\thesection}{1em}{}
\titleformat{\subsection}[hang]{\large\sffamily\color{nyuDarkPurple}}{\thesubsection}{1em}{}
\titleformat{\subsubsection}[hang]{\normalsize\sffamily\color{gray}}{\thesubsubsection}{1em}{}

\setcounter{secnumdepth}{0}

% math commands
\newcommand{\R}{\mathbb{R}}
\newcommand{\E}{\mathbb{E}}
\DeclareMathOperator{\Var}{Var}
\newcommand{\bs}[1]{\boldsymbol{#1}}
\newcommand{\bv}[1]{\mathbf{#1}}

\begin{document}

\begin{center}
	\normalsize
	New York University
	\medskip

	\large
	CS-GY 6763: Homework 2

	Due Wednesday, February 25, 2025, 11:59pm ET.
	\medskip

	\normalsize
	\noindent \emph{Collaboration is allowed on this problem set, but solutions must be written up individually. Please list collaborators for each problem separately, or write ``No Collaborators'' if you worked alone.}
\end{center}

\subsection{Problem 1: Exploring Concentration Bounds}
\textbf{(8 pts)} Let $X$ be a random variable uniformly distributed on $[0,1]$. Since we know $X$'s distribution exactly, we can easily check that $\Pr[X \geq 7/8] = 1/8$. But let's take a look at what various concentration inequalities would predict about this probability using less information about $X$.

\begin{enumerate}[label=\arabic*.]

	\item Give an upper bound on $\Pr[X \geq 7/8]$ using Markov's inequality.

\spara{Pr. 1.1}

\[
X\sim \text{Uniform}[0,1],
\qquad
\E[X]=0.5,
\qquad
a=\frac{7}{8}.
\]

Markov:
\[
\Pr(X\ge a)\le \frac{\E[X]}{a}.
\]

\[
\Pr\!\left(X\ge \frac{7}{8}\right)
\le
\frac{\E[X]}{7/8}
=
\frac{1/2}{7/8}
=
\frac{1}{2}\cdot \frac{8}{7}
=
\frac{4}{7}
\approx 0.571.
\]

\[
\boxed{\Pr[X\ge 7/8]\le \frac{4}{7}.}
\]


	\item Give an upper bound on $\Pr[X \geq 7/8]$ using Chebyshev's inequality.

\spara{Pr. 1.2}

Chebyshev:
\[
\Pr(|X-\E[X]|\ge t)\le \frac{\sigma^2}{t^2},
\qquad
\sigma^2=\Var(X).
\]

For $X\sim \text{Uniform}[0,1]$,
\[
\Var(X)=\frac{1}{12},
\qquad
\E[X]=\frac{1}{2}.
\]

We want $\Pr(X\ge 7/8)$. Note
\[
X\ge \frac{7}{8}
\ \Rightarrow\
X-\E[X]\ge \frac{7}{8}-\frac{1}{2}=\frac{3}{8}
\ \Rightarrow\
|X-\E[X]|\ge \frac{3}{8}.
\]
So
\[
\Pr\!\left(X\ge \frac{7}{8}\right)
\le
\Pr\!\left(|X-\E[X]|\ge \frac{3}{8}\right)
\le
\frac{1/12}{(3/8)^2}
=
\frac{1/12}{9/64}
=
\frac{64}{108}
=
\frac{16}{27}
\approx 0.593.
\]

\[
\boxed{\Pr[X\ge 7/8]\le \frac{16}{27}.}
\]


	\item Give an upper bound on $\Pr[X \geq 7/8]$ by applying Markov's inequality to the random variable $X^2$ (the ``raw'' second moment). Note that this is slightly different than using Chebyshev's inequality, which applies Markov to the ``central'' second moment $(X - \E[X])^2$.

\spara{Pr. 1.3}

Since $X\ge 0$,
\[
\left\{X\ge \frac{7}{8}\right\}
=
\left\{X^2 \ge \left(\frac{7}{8}\right)^2\right\}
=
\left\{X^2 \ge \frac{49}{64}\right\}.
\]

Markov on $X^2$:
\[
\Pr(X^2\ge a^2)\le \frac{\E[X^2]}{a^2},
\qquad
a=\frac{7}{8}.
\]

Compute the raw second moment:
\[
\E[X^2]=\int_0^1 x^2\,dx
=
\left[\frac{x^3}{3}\right]_0^1
=
\frac{1}{3}.
\]

Therefore
\[
\Pr\!\left(X\ge \frac{7}{8}\right)
=
\Pr\!\left(X^2\ge \frac{49}{64}\right)
\le
\frac{1/3}{49/64}
=
\frac{64}{147}
\approx 0.435.
\]

\[
\boxed{\Pr[X\ge 7/8]\le \frac{64}{147}.}
\]


	\item What happens for higher moments? Apply Markov's inequality to $X^q$ for $q = 3,4,\ldots,10$. Describe what you see in a sentence. What value of $q$ gives the tightest bound?

\spara{Pr. 1.4}

For $q\ge 1$,
\[
\Pr\!\left(X\ge \frac{7}{8}\right)
=
\Pr\!\left(X^q \ge \left(\frac{7}{8}\right)^q\right)
\le
\frac{\E[X^q]}{(7/8)^q}.
\]

For uniform $[0,1]$,
\[
\E[X^q]=\int_0^1 x^q\,dx=\frac{1}{q+1}.
\]

So the bound is
\[
\Pr\!\left(X\ge \frac{7}{8}\right)\le
\frac{1}{(q+1)(7/8)^q}.
\]

Values (table from my paper):
\[
\begin{array}{c|c}
q & \text{Markov on }X^q \\
\hline
3 & 0.373\\
4 & 0.341\\
5 & 0.324\\
6 & 0.318\\
7 & 0.318\\
8 & 0.323\\
9 & 0.332\\
10 & 0.340
\end{array}
\]

\[
\boxed{\text{Minimum is at } q=6 \text{ (ties with } q=7\text{).}}
\]

Sentence:
\[
\text{The bound improves up to about }q=6\text{ then gets worse because }(8/7)^q\text{ grows.}
\]


	\item One takeaway here is that, depending on the random variable being studied, it is not always optimal to use the variance as a deviation measure to show concentration. Markov's inequality can be used with any monotone function $g$, and as we see above, different choices might give better bounds. Exhibit a monotone function $g$ so that applying Markov's inequality to $g(X)$ gives as tight an upper bound on $\Pr[X \geq 7/8]$ as you can. Maximum points if you can get $\Pr[X \geq 7/8] \leq 1/8$, which would be the best possible.

\spara{Pr. 1.5}

For monotone $g$:
\[
\Pr(X\ge a)=\Pr(g(X)\ge g(a))\le \frac{\E[g(X)]}{g(a)}.
\]

Pick
\[
g(x)=
\begin{cases}
1 & x\ge 7/8\\
0 & \text{else}
\end{cases}
\qquad\Rightarrow\qquad
g \text{ monotone.}
\]

Then $g(X)\sim \text{Bernoulli}(p)$ where $p=\Pr(X\ge 7/8)=1/8$.
So
\[
\E[g(X)] = 1\cdot p + 0\cdot (1-p) = \frac{1}{8}.
\]

Markov (threshold $g(7/8)=1$):
\[
\Pr(X\ge 7/8)=\Pr(g(X)\ge 1)\le \frac{\E[g(X)]}{1}=\frac{1}{8}.
\]

\[
\boxed{\Pr[X\ge 7/8]\le \frac{1}{8}.}
\]

\spara{Notes for Problem 1}
\begin{itemize}
\item Copied from handwritten pages 1--3.
\item The numeric table in part (4) is copied as written, so small rounding differences possible.
\end{itemize}

\end{enumerate}



\subsection{Problem 2: Boosting Success Probability}
\textbf{(12 pts)}

\spara{Pr. 2}

Setup:
\[
\Pr(|\mathcal{A}[X]-f(X)|\le \epsilon)\ge \frac{2}{3}
\quad\Rightarrow\quad
\Pr(|\mathcal{A}[X]-f(X)|>\epsilon)\le \frac{1}{3}.
\]

Run $\mathcal{A}$ independently $r$ times.

Let
\[
\bar{S}=\#\text{ of failures},\qquad S=\#\text{ of successes},\qquad S+\bar{S}=r.
\]

Median is ``bad'' only if at least half are bad:
\[
\{|M-f(X)|>\epsilon\}\ \Rightarrow\ \left\{\bar{S}\ge \frac{r}{2}\right\}.
\]

Also
\[
\E[\bar{S}] \le \frac{r}{3}.
\]

Chernoff (one-sided upper tail):
\[
\Pr(\bar{S}\ge (1+\eta)\mu)\le \exp\!\left(-\frac{\eta^2\mu}{2+\eta}\right),
\qquad
\mu=\E[\bar{S}].
\]

We want $\bar{S}\ge r/2$ with $\mu=r/3$ (worst case). Solve:
\[
\frac{r}{2}=(1+\eta)\frac{r}{3}
\Rightarrow
1+\eta=\frac{3}{2}
\Rightarrow
\eta=\frac{1}{2}.
\]

So
\[
\Pr\!\left(\bar{S}\ge \frac{r}{2}\right)
\le
\exp\!\left(-\frac{(1/2)^2\cdot (r/3)}{2+1/2}\right)
=
\exp\!\left(-\frac{r}{30}\right).
\]

Set $\exp(-r/30)\le \delta$:
\[
e^{-r/30}\le \delta
\Rightarrow
r\ge 30\ln(1/\delta).
\]

Therefore
\[
\Pr(|M-f(X)|\le \epsilon)\ge 1-\delta,
\qquad
\text{with } r=O(\log(1/\delta)).
\qquad \Box
\]

\spara{Notes for Problem 2}
\begin{itemize}
\item Copied from handwritten pages 4--5.
\item I used $\bar{S}$ for failures like on my paper.
\end{itemize}



\subsection{Problem 3: Hashing around the clock}
\textbf{(15 pts)}

\begin{enumerate}[label=\arabic*.]

\item Suppose we have $n$ servers initially. When a new server is added to the system, what is the expected number of data items that need to be relocated?

\spara{Pr. 3.1}

\textbf{Circle picture (words).}
Put $[0,1)$ on a circle. Hash every server to a point. Hash every item to a point.
Each item goes to the first server clockwise.

When we add a new server, only items in the arc immediately before the new server move to it.

Let $I_j$ be indicator that item $j$ moves:
\[
I_j =
\begin{cases}
1 & \text{item $j$ gets assigned to the new server}\\
0 & \text{otherwise.}
\end{cases}
\qquad
\#\text{moved}=\sum_{j=1}^m I_j.
\]

After insertion there are $(n+1)$ servers total. By symmetry,
\[
\Pr(I_j=1)=\frac{1}{n+1}
\quad\Rightarrow\quad
\E[I_j]=\frac{1}{n+1}.
\]
So
\[
\E[\#\text{moved}] = \sum_{j=1}^m \E[I_j]=\sum_{j=1}^m \frac{1}{n+1}=\frac{m}{n+1}.
\]
\[
\boxed{\E[\#\text{items relocated}] = \frac{m}{n+1}.}
\]


\item Show that, with probability $> 9/10$, no server ``owns'' more than an $O(\log n / n)$ fraction of the interval $[0,1]$.

\spara{Pr. 3.2}

Let $G_{\max}$ be the maximum gap between consecutive server points on the circle.

Fix a server. Probability that the next server clockwise is at distance at least $t$ is
\[
\Pr(\text{gap}\ge t)=(1-t)^{n-1}
\]
(since none of the other $(n-1)$ servers fall in that length-$t$ interval).

Union bound over $n$ gaps:
\[
\Pr(G_{\max}\ge t)\le n(1-t)^{n-1}\le n e^{-(n-1)t}.
\]

Pick $t=c\frac{\log n}{n}$:
\[
\Pr\!\left(G_{\max}\ge c\frac{\log n}{n}\right)
\le n\exp\!\left(-(n-1)c\frac{\log n}{n}\right)
\le n^{1-c/2}.
\]
Choose $c=6$, then $n^{1-c/2}=n^{-2}<1/10$ for $n\ge 4$.

So w.p.\ $>9/10$,
\[
G_{\max}\le 6\frac{\log n}{n}.
\]
\[
\boxed{\text{No server owns more than }O(\log n/n)\text{ w.p. } >9/10.}
\]


\item Show that if we have $n$ servers and $m$ items with $m > n$, the maximum load on any server is no more than $O\!\left(\frac{m}{n}\log n\right)$ with probability $> 9/10$.

\spara{Pr. 3.3}

Let $L_i$ be the interval length owned by server $i$.
Conditioned on the server locations, for fixed $L_i$,
\[
\text{Load}_i \sim \mathrm{Binomial}(m,L_i),
\qquad
\E[\text{Load}_i\mid L_i]=mL_i.
\]

Let $t=6\frac{\log n}{n}$ and event
\[
E:=\left\{\max_i L_i\le t\right\}.
\]
From Pr.\ 3.2, $\Pr(E)>9/10$ (for $n$ not tiny).

On $E$, $\mu:=mt=6\frac{m}{n}\log n$ bounds all means.

Use Chernoff-style upper tail:
\[
\Pr(\text{Load}_i\ge 2\mu\mid E)\le \exp(-\mu/3).
\]
Since $m>n$, $\mu\ge 6\log n$, so $\exp(-\mu/3)\le n^{-2}$.
Union bound:
\[
\Pr\left(\exists i:\ \text{Load}_i\ge 2\mu\ \middle|\ E\right)\le n\cdot n^{-2}=\frac{1}{n}.
\]
So (conditioned on $E$) w.p.\ $\ge 1-1/n$,
\[
\max_i \text{Load}_i\le 2\mu = 12\frac{m}{n}\log n.
\]

Therefore overall with prob $>9/10$,
\[
\boxed{\max_i \text{Load}_i = O\!\left(\frac{m}{n}\log n\right).}
\]

\spara{Notes for Problem 3}
\begin{itemize}
\item This one was not in my handwritten pages, so I wrote it cleanly.
\end{itemize}

\end{enumerate}



\subsection{Problem 4: Johnson--Lindenstrauss Approximates Inner Products}
\textbf{(10 pts)}

\begin{enumerate}[label=\arabic*.]

\item Suppose that $\Pi$ is a Johnson--Lindenstrauss matrix with $k = O\!\left(\frac{\log(1/\delta)}{\epsilon^2}\right)$ rows. Prove that for any $x,y$,
\[
\big|\langle x,y\rangle - \langle \Pi x,\Pi y\rangle\big|
\leq \epsilon \left(\|x\|_2^2 + \|y\|_2^2\right)
\]
with probability at least $1-\delta$.

\spara{Pr. 4.1}

Lemma (Distributional JL Lemma): choose $\Pi$ so each entry is $\frac{1}{\sqrt{k}}N(0,1)$.
If $k=O\!\left(\frac{\log(1/\delta)}{\epsilon^2}\right)$ then w.p.\ $(1-\delta)$,
\[
(1-\epsilon)\|v\|_2^2 \le \|\Pi v\|_2^2 \le (1+\epsilon)\|v\|_2^2.
\]

Product vs Norm (polarization):
\[
\langle x,y\rangle = \frac{\|x+y\|_2^2-\|x-y\|_2^2}{4},
\qquad
\langle \Pi x,\Pi y\rangle = \frac{\|\Pi(x+y)\|_2^2-\|\Pi(x-y)\|_2^2}{4}.
\]

So
\[
|\langle x,y\rangle - \langle \Pi x,\Pi y\rangle|
=
\frac{1}{4}\left|
(\|x+y\|_2^2-\|x-y\|_2^2)
-
(\|\Pi(x+y)\|_2^2-\|\Pi(x-y)\|_2^2)
\right|.
\]

Triangle inequality:
\[
\le
\frac{1}{4}\left(
\big|\|x+y\|_2^2-\|\Pi(x+y)\|_2^2\big|
+
\big|\|x-y\|_2^2-\|\Pi(x-y)\|_2^2\big|
\right).
\]

If JL holds for $(x+y)$ and $(x-y)$ (use $\delta/2$ each then union bound), then
\[
\big|\|x\pm y\|_2^2-\|\Pi(x\pm y)\|_2^2\big|\le \epsilon\|x\pm y\|_2^2.
\]
So
\[
|\langle x,y\rangle - \langle \Pi x,\Pi y\rangle|
\le
\frac{\epsilon}{4}\left(\|x+y\|_2^2+\|x-y\|_2^2\right).
\]

Parallelogram identity:
\[
\|x+y\|_2^2+\|x-y\|_2^2 = 2(\|x\|_2^2+\|y\|_2^2).
\]
Hence
\[
|\langle x,y\rangle - \langle \Pi x,\Pi y\rangle|
\le
\frac{\epsilon}{2}(\|x\|_2^2+\|y\|_2^2).
\]
(Absorb the constant by redefining $\epsilon$.) \(\Box\)

\spara{Notes (Pr. 4.1)}
\begin{itemize}
\item Copied from handwritten pages 6--7.
\item I used an $\epsilon_{\text{new}}$ idea on paper to absorb constants.
\end{itemize}


\item Show that a better bound can be obtained by augmenting our sketches with the norms of each vector.

\spara{Pr. 4.2}

Sketch:
\[
\text{sketch}(x) = (\Pi x,\ \|x\|_2),
\qquad
\text{sketch}(y) = (\Pi y,\ \|y\|_2).
\]

Normalize:
\[
x_n=\frac{x}{\|x\|_2},
\qquad
y_n=\frac{y}{\|y\|_2},
\qquad
\|x_n\|_2=\|y_n\|_2=1.
\]
Then
\[
\langle x_n,y_n\rangle = \frac{\langle x,y\rangle}{\|x\|_2\|y\|_2},
\qquad
\Pi x_n = \frac{\Pi x}{\|x\|_2},
\qquad
\Pi y_n = \frac{\Pi y}{\|y\|_2}.
\]

Using Pr.\ 4.1 on $(x_n,y_n)$ (and absorbing constants into $\epsilon$):
\[
|\langle x_n,y_n\rangle - \langle \Pi x_n,\Pi y_n\rangle|
\le \epsilon.
\]

Define
\[
c=\langle \Pi x_n,\Pi y_n\rangle,
\qquad
z := \|x\|_2\|y\|_2\,c.
\]

Then
\[
|\langle x,y\rangle - z|
=
\|x\|_2\|y\|_2\cdot |\langle x_n,y_n\rangle - c|
\le
\epsilon\|x\|_2\|y\|_2.
\]
\(\Box\)

\spara{Notes (Pr. 4.2)}
\begin{itemize}
\item Copied from handwritten pages 8--9.
\end{itemize}

\end{enumerate}



\subsection{Problem 5: Analyzing Sign-JL and JL for Inner Products}
\textbf{(15 pts)}

\begin{enumerate}[label=\arabic*.]

\item Prove that for any vector $\bv{x}\in \R^n$, $\E[\|\bs{\Pi}\bv{x}\|_2^2] = \|\bv{x}\|_2^2$ and that
\[
\Var\!\left[\|\bs{\Pi}\bv{x}\|_2^2\right] \leq \frac{2}{k}\|\bv{x}\|_2^4.
\]

\spara{Pr. 5.1}

Each entry is $\pm \frac{1}{\sqrt{k}}$:
\[
\Pr(\Pi_{ij}=+1/\sqrt{k})=\frac{1}{2},
\qquad
\Pr(\Pi_{ij}=-1/\sqrt{k})=\frac{1}{2}.
\]
So
\[
\E[\Pi_{ij}]=0,
\qquad
\E[\Pi_{ij}^2]=\frac{1}{k}.
\]

Let $(\Pi x)_i$ be row $i$ dot-product:
\[
(\Pi x)_i=\sum_{j=1}^n \Pi_{ij}x_j.
\]
Then
\[
\|\Pi x\|_2^2=\sum_{i=1}^k (\Pi x)_i^2.
\]

Use $\Pi_{ij}=\frac{1}{\sqrt{k}}O_{ij}$ where $O_{ij}\in\{\pm 1\}$ i.i.d.
Define
\[
Y_i := \left(\sum_{j=1}^n O_{ij}x_j\right)^2,
\qquad
Z:=\|\Pi x\|_2^2=\frac{1}{k}\sum_{i=1}^k Y_i.
\]

\textbf{Mean.}
Expand $Y_i$:
\[
Y_i
=
\left(\sum_{j=1}^n O_{ij}x_j\right)^2
=
\sum_{j=1}^n O_{ij}^2 x_j^2
+
\sum_{j\ne \ell}O_{ij}O_{i\ell}x_jx_\ell.
\]
Take expectation:
\[
\E[O_{ij}^2]=1,
\qquad
\E[O_{ij}O_{i\ell}]=\E[O_{ij}]\E[O_{i\ell}]=0\quad (j\ne \ell),
\]
so
\[
\E[Y_i]=\sum_{j=1}^n x_j^2=\|x\|_2^2.
\]
Then
\[
\E[Z]=\E\!\left[\frac{1}{k}\sum_{i=1}^k Y_i\right]
=
\frac{1}{k}\sum_{i=1}^k \E[Y_i]
=
\|x\|_2^2.
\]
So
\[
\boxed{\E[\|\Pi x\|_2^2]=\|x\|_2^2.}
\]

\textbf{Variance.}
\[
\Var(Z)=\E[Z^2]-\E[Z]^2,
\qquad
Z=\frac{1}{k}\sum_{i=1}^k Y_i.
\]

Square:
\[
Z^2
=
\left(\frac{1}{k}\sum_{i=1}^k Y_i\right)^2
=
\frac{1}{k^2}\left(\sum_{i=1}^k Y_i^2 + 2\sum_{i<j}Y_iY_j\right).
\]

Using independence of rows (as on my notes):
\[
\E[Y_iY_j]=\E[Y_i]\E[Y_j]\quad (i\ne j),
\]
so the cross term cancels in $\Var(Z)$ and we get the reduction
\[
\Var(Z)=\frac{1}{k}\Var(Y_1).
\]

Now use the 4th-moment expansion:
\[
\E[Y_1^2]
=
\E\!\left[\left(\sum_{j=1}^n O_{1j}x_j\right)^4\right]
=
\sum_p x_p^4 + 6\sum_{p<q}x_p^2x_q^2
\le 3\|x\|_2^4.
\]
Also $\E[Y_1]=\|x\|_2^2$, so
\[
\Var(Y_1)=\E[Y_1^2]-\E[Y_1]^2 \le 3\|x\|_2^4-\|x\|_2^4=2\|x\|_2^4.
\]
Thus
\[
\Var(Z)=\frac{1}{k}\Var(Y_1)\le \frac{2}{k}\|x\|_2^4.
\]

\[
\boxed{\Var(\|\Pi x\|_2^2)\le \frac{2}{k}\|x\|_2^4.}
\]

\spara{Notes (Pr. 5.1)}
\begin{itemize}
\item Copied from handwritten pages 10--14. The key split was writing $Z=\|\Pi x\|_2^2=\frac{1}{k}\sum_i Y_i$.
\item The 4th moment line with the ``$6$'' pairing term was the main hard step in my notes.
\item Some intermediate circled terms on my paper were hard to read; I kept the same structure and then used the standard cancellation to get $\Var(Z)=\frac{1}{k}\Var(Y_1)$.
\end{itemize}


\item Use the above to prove that
\[
\Pr\!\left[\left|\|\bs{\Pi}\bv{x}\|_2^2 - \|\bv{x}\|_2^2\right| \geq \epsilon\|\bv{x}\|_2^2 \right] \leq \delta
\]
as long as we choose $k = O\!\left(\frac{1}{\delta\epsilon^2}\right)$.

\spara{Pr. 5.2}

Let
\[
Z=\|\Pi x\|_2^2,
\qquad
\E[Z]=\|x\|_2^2,
\qquad
\Var(Z)\le \frac{2}{k}\|x\|_2^4.
\]

Chebyshev:
\[
\Pr(|Z-\E[Z]|\ge t)\le \frac{\Var(Z)}{t^2}.
\]

Set
\[
t=\epsilon\|x\|_2^2.
\]
Then
\[
\Pr\!\left(|\|\Pi x\|_2^2-\|x\|_2^2|\ge \epsilon\|x\|_2^2\right)
\le
\frac{(2/k)\|x\|_2^4}{\epsilon^2\|x\|_2^4}
=
\frac{2}{k\epsilon^2}.
\]

Make it $\le \delta$:
\[
\frac{2}{k\epsilon^2}\le \delta
\Rightarrow
k\ge \frac{2}{\delta\epsilon^2}
\Rightarrow
k=O\!\left(\frac{1}{\delta\epsilon^2}\right).
\]
\(\Box\)

\spara{Notes (Pr. 5.2)}
\begin{itemize}
\item Copied from handwritten page 15 (Chebyshev plug-in).
\end{itemize}

\end{enumerate}

\end{document}